% --- Archon physics formalization source begin ---
% archon:physics
% archon:covers PhyXMiniProblems/problem_phyx_mini_0839.lean
% archon:source-report reports/phyx_mini/problem_phyx_mini_0839.source.json

\chapter{Physics problem phyx\_mini\_0839}
\label{ch:PhyXMiniProblems_problem_phyx_mini_0839}

\paragraph{Problem source.}
\#\# Physical scenario

Your physics assignment is to figure out a way to use electricity to launch a small \textbackslash{}( 6.0\textbackslash{}\textbackslash{},\textbackslash{}\textbackslash{}mathrm\{cm\} \textbackslash{})-long plastic drink stirrer. You decide that you'll charge the little plastic rod by rubbing it with fur, then hold it near a long, charged wire, as shown in figure. When you let go, the electric force of the wire on the plastic rod will shoot it away. Suppose you can uniformly charge the plastic stirrer to \textbackslash{}( 10\textbackslash{}\textbackslash{},\textbackslash{}\textbackslash{}mathrm\{nC\} \textbackslash{}) and that the linear charge density of the long wire is \textbackslash{}( 1.0 	imes 10\textasciicircum{}\{-7\}\textbackslash{}\textbackslash{},\textbackslash{}\textbackslash{}mathrm\{C/m\} \textbackslash{}).

\#\# Question

What is the net electric force on the plastic stirrer if the end closest to the wire is \textbackslash{}( 2.0\textbackslash{}\textbackslash{},\textbackslash{}\textbackslash{}mathrm\{cm\} \textbackslash{}) away?

\#\# Answer choices

- A: A: 1.3 \textbackslash{}times 10\textasciicircum{}\{-5\}N

- B: B: 5.3 \textbackslash{}times 10\textasciicircum{}\{-4\}N

- C: C: 1.0 \textbackslash{}times 10\textasciicircum{}\{-5\}N

- D: D: 4.2 \textbackslash{}times 10\textasciicircum{}\{-4\}N

\#\# Figure caption (auxiliary; use the image as primary evidence)

The image shows two main objects: a charged vertical surface and a horizontal plastic stirrer.

1. **Charged Surface**: 
   - Positioned on the left, it has red plus signs indicating positive charge.
   - There is an annotation pointing to it with the charge density: \textbackslash{}( \textbackslash{}lambda = 1.0 \textbackslash{}times 10\textasciicircum{}\{-7\} \textbackslash{}, \textbackslash{}text\{C/m\} \textbackslash{}).
   
2. **Plastic Stirrer**:
   - Placed horizontally to the right of the charged surface.
   - Labeled with the text "Plastic stirrer."

3. **Distances**: 
   - The distance between the charged surface and the plastic stirrer is marked as 2.0 cm.
   - The length from the left edge of the stirrer is marked as 6.0 cm.

\#\# Dataset metadata

Electrostatics; Physical Model Grounding Reasoning, Multi-Formula Reasoning

\paragraph{Recorded answer/context.}
D: D: 4.2 \textbackslash{}times 10\textasciicircum{}\{-4\}N

\paragraph{Figure/image path.}
/root/proposal\_for\_physic/science-mango/phyx\_mini\_run/phyx\_data/test\_image/839.png

\paragraph{Formalization target.}
create a compiling Lean file with sorry bodies at `PhyXMiniProblems/problem\_phyx\_mini\_0839.lean`. The Lean declarations must preserve the physical quantities, dimensions or dimensional roles, figure labels, governing-law hypotheses, and final relation expressed by this problem.
Use Mathlib/Physlib names found through LeanExplore where available. If a physics API is missing, introduce faithful local abstractions rather than scalar placeholder aliases.

\begin{theorem}[Physics formalization target]
\label{thm:physics:phyx_mini_0839:target}
\lean{PhyXMiniProblems.ProblemPhyXMini0839.netElectricForce_isAwayFromWire_and_agreesWithChoiceD}
\uses{def:physics:phyx-mini-0839:phyxminiproblems-problemphyxmini0839-forceinnewtons, def:physics:phyx-mini-0839:phyxminiproblems-problemphyxmini0839-chargedstirrersetup, def:physics:phyx-mini-0839:phyxminiproblems-problemphyxmini0839-haschargedstirrerproblemdata, def:physics:phyx-mini-0839:phyxminiproblems-problemphyxmini0839-matcheschargedstirrerscenario, def:physics:phyx-mini-0839:phyxminiproblems-problemphyxmini0839-matcheschargedstirrerprimaryfigure, def:physics:phyx-mini-0839:phyxminiproblems-problemphyxmini0839-satisfiesinfinitelinechargeelectrostatics, def:physics:phyx-mini-0839:phyxminiproblems-problemphyxmini0839-answerchoiceforceinnewtons, def:physics:phyx-mini-0839:phyxminiproblems-problemphyxmini0839-choiceddisplaytoleranceinnewtons}
The assigned autoformalize agent should translate this physics problem into Lean declarations in the covered file, with theorem and lemma proof bodies written as `by sorry`.
\end{theorem}
\begin{proof}
This is an autoformalization task, not a proof task. Produce statements that can later be proved by the physics prover without weakening the physical model.
\end{proof}

% --- Archon declaration topology begin ---
\section{Declaration topology}
\begin{definition}[force Dimension]
  \label{def:physics:phyx-mini-0839:phyxminiproblems-problemphyxmini0839-forcedimension}
  \lean{PhyXMiniProblems.ProblemPhyXMini0839.forceDimension}
  Dimension of force, M L T\textasciicircum{}-2
\end{definition}
\begin{proof}
  This is the declared model definition of force Dimension.
\end{proof}
\begin{definition}[line Charge Density Dimension]
  \label{def:physics:phyx-mini-0839:phyxminiproblems-problemphyxmini0839-linechargedensitydimension}
  \lean{PhyXMiniProblems.ProblemPhyXMini0839.lineChargeDensityDimension}
  Dimension of linear charge density, C L\textasciicircum{}-1
\end{definition}
\begin{proof}
  This is the declared model definition of line Charge Density Dimension.
\end{proof}
\begin{definition}[electric Field Dimension]
  \label{def:physics:phyx-mini-0839:phyxminiproblems-problemphyxmini0839-electricfielddimension}
  \lean{PhyXMiniProblems.ProblemPhyXMini0839.electricFieldDimension}
  \uses{def:physics:phyx-mini-0839:phyxminiproblems-problemphyxmini0839-forcedimension}
  Dimension of electric-field magnitude, N/C
\end{definition}
\begin{proof}
  This is the declared model definition of electric Field Dimension.
\end{proof}
\begin{definition}[coulomb Constant Dimension]
  \label{def:physics:phyx-mini-0839:phyxminiproblems-problemphyxmini0839-coulombconstantdimension}
  \lean{PhyXMiniProblems.ProblemPhyXMini0839.coulombConstantDimension}
  \uses{def:physics:phyx-mini-0839:phyxminiproblems-problemphyxmini0839-forcedimension}
  Dimension of Coulomb's constant, N m\textasciicircum{}2 / C\textasciicircum{}2
\end{definition}
\begin{proof}
  This is the declared model definition of coulomb Constant Dimension.
\end{proof}
\begin{definition}[Length Quantity]
  \label{def:physics:phyx-mini-0839:phyxminiproblems-problemphyxmini0839-lengthquantity}
  \lean{PhyXMiniProblems.ProblemPhyXMini0839.LengthQuantity}
  A nonnegative, unit-independent physical length
\end{definition}
\begin{proof}
  This is the declared model definition of Length Quantity.
\end{proof}
\begin{definition}[Charge Magnitude Quantity]
  \label{def:physics:phyx-mini-0839:phyxminiproblems-problemphyxmini0839-chargemagnitudequantity}
  \lean{PhyXMiniProblems.ProblemPhyXMini0839.ChargeMagnitudeQuantity}
  A nonnegative, unit-independent electric-charge magnitude
\end{definition}
\begin{proof}
  This is the declared model definition of Charge Magnitude Quantity.
\end{proof}
\begin{definition}[Line Charge Density Magnitude Quantity]
  \label{def:physics:phyx-mini-0839:phyxminiproblems-problemphyxmini0839-linechargedensitymagnitudequantity}
  \lean{PhyXMiniProblems.ProblemPhyXMini0839.LineChargeDensityMagnitudeQuantity}
  \uses{def:physics:phyx-mini-0839:phyxminiproblems-problemphyxmini0839-linechargedensitydimension}
  A nonnegative linear-charge-density magnitude
\end{definition}
\begin{proof}
  This is the declared model definition of Line Charge Density Magnitude Quantity.
\end{proof}
\begin{definition}[Force Magnitude Quantity]
  \label{def:physics:phyx-mini-0839:phyxminiproblems-problemphyxmini0839-forcemagnitudequantity}
  \lean{PhyXMiniProblems.ProblemPhyXMini0839.ForceMagnitudeQuantity}
  \uses{def:physics:phyx-mini-0839:phyxminiproblems-problemphyxmini0839-forcedimension}
  A nonnegative physical force magnitude
\end{definition}
\begin{proof}
  This is the declared model definition of Force Magnitude Quantity.
\end{proof}
\begin{definition}[Coulomb Constant Quantity]
  \label{def:physics:phyx-mini-0839:phyxminiproblems-problemphyxmini0839-coulombconstantquantity}
  \lean{PhyXMiniProblems.ProblemPhyXMini0839.CoulombConstantQuantity}
  \uses{def:physics:phyx-mini-0839:phyxminiproblems-problemphyxmini0839-coulombconstantdimension}
  A nonnegative value of Coulomb's constant with its physical dimension
\end{definition}
\begin{proof}
  This is the declared model definition of Coulomb Constant Quantity.
\end{proof}
\begin{definition}[length In Meters]
  \label{def:physics:phyx-mini-0839:phyxminiproblems-problemphyxmini0839-lengthinmeters}
  \lean{PhyXMiniProblems.ProblemPhyXMini0839.lengthInMeters}
  \uses{def:physics:phyx-mini-0839:phyxminiproblems-problemphyxmini0839-lengthquantity}
  Coherent-SI metre readout of a physical length
\end{definition}
\begin{proof}
  This is the declared model definition of length In Meters.
\end{proof}
\begin{definition}[charge In Coulombs]
  \label{def:physics:phyx-mini-0839:phyxminiproblems-problemphyxmini0839-chargeincoulombs}
  \lean{PhyXMiniProblems.ProblemPhyXMini0839.chargeInCoulombs}
  \uses{def:physics:phyx-mini-0839:phyxminiproblems-problemphyxmini0839-chargemagnitudequantity}
  Coherent-SI coulomb readout of a charge magnitude
\end{definition}
\begin{proof}
  This is the declared model definition of charge In Coulombs.
\end{proof}
\begin{definition}[line Charge Density In Coulombs Per Meter]
  \label{def:physics:phyx-mini-0839:phyxminiproblems-problemphyxmini0839-linechargedensityincoulombspermeter}
  \lean{PhyXMiniProblems.ProblemPhyXMini0839.lineChargeDensityInCoulombsPerMeter}
  \uses{def:physics:phyx-mini-0839:phyxminiproblems-problemphyxmini0839-linechargedensitymagnitudequantity}
  Coherent-SI coulomb-per-metre readout of a linear charge density
\end{definition}
\begin{proof}
  This is the declared model definition of line Charge Density In Coulombs Per Meter.
\end{proof}
\begin{definition}[force In Newtons]
  \label{def:physics:phyx-mini-0839:phyxminiproblems-problemphyxmini0839-forceinnewtons}
  \lean{PhyXMiniProblems.ProblemPhyXMini0839.forceInNewtons}
  \uses{def:physics:phyx-mini-0839:phyxminiproblems-problemphyxmini0839-forcemagnitudequantity}
  Coherent-SI newton readout of a physical force magnitude
\end{definition}
\begin{proof}
  This is the declared model definition of force In Newtons.
\end{proof}
\begin{definition}[coulomb Constant In Newton Meters Squared Per Coulomb Squared]
  \label{def:physics:phyx-mini-0839:phyxminiproblems-problemphyxmini0839-coulombconstantinnewtonmeterssquaredpercoulombsquared}
  \lean{PhyXMiniProblems.ProblemPhyXMini0839.coulombConstantInNewtonMetersSquaredPerCoulombSquared}
  \uses{def:physics:phyx-mini-0839:phyxminiproblems-problemphyxmini0839-coulombconstantquantity}
  Coherent-SI readout of Coulomb's constant in N m\textasciicircum{}2 / C\textasciicircum{}2
\end{definition}
\begin{proof}
  This is the declared model definition of coulomb Constant In Newton Meters Squared Per Coulomb Squared.
\end{proof}
\begin{definition}[Figure Object]
  \label{def:physics:phyx-mini-0839:phyxminiproblems-problemphyxmini0839-figureobject}
  \lean{PhyXMiniProblems.ProblemPhyXMini0839.FigureObject}
  The two physical objects shown in the primary image
\end{definition}
\begin{proof}
  This is the declared model definition of Figure Object.
\end{proof}
\begin{definition}[Figure Label]
  \label{def:physics:phyx-mini-0839:phyxminiproblems-problemphyxmini0839-figurelabel}
  \lean{PhyXMiniProblems.ProblemPhyXMini0839.FigureLabel}
  Literal annotations visible in image 839.png
\end{definition}
\begin{proof}
  This is the declared model definition of Figure Label.
\end{proof}
\begin{definition}[Charge Sign]
  \label{def:physics:phyx-mini-0839:phyxminiproblems-problemphyxmini0839-chargesign}
  \lean{PhyXMiniProblems.ProblemPhyXMini0839.ChargeSign}
  Sign of an object's electric charge
\end{definition}
\begin{proof}
  This is the declared model definition of Charge Sign.
\end{proof}
\begin{definition}[Wire Model]
  \label{def:physics:phyx-mini-0839:phyxminiproblems-problemphyxmini0839-wiremodel}
  \lean{PhyXMiniProblems.ProblemPhyXMini0839.WireModel}
  Idealization of the long charged object
\end{definition}
\begin{proof}
  This is the declared model definition of Wire Model.
\end{proof}
\begin{definition}[Stirrer Charge Distribution]
  \label{def:physics:phyx-mini-0839:phyxminiproblems-problemphyxmini0839-stirrerchargedistribution}
  \lean{PhyXMiniProblems.ProblemPhyXMini0839.StirrerChargeDistribution}
  Distribution of charge along the plastic stirrer
\end{definition}
\begin{proof}
  This is the declared model definition of Stirrer Charge Distribution.
\end{proof}
\begin{definition}[Wire Stirrer Orientation]
  \label{def:physics:phyx-mini-0839:phyxminiproblems-problemphyxmini0839-wirestirrerorientation}
  \lean{PhyXMiniProblems.ProblemPhyXMini0839.WireStirrerOrientation}
  Relative orientation of the wire and the stirrer
\end{definition}
\begin{proof}
  This is the declared model definition of Wire Stirrer Orientation.
\end{proof}
\begin{definition}[Radial Force Direction]
  \label{def:physics:phyx-mini-0839:phyxminiproblems-problemphyxmini0839-radialforcedirection}
  \lean{PhyXMiniProblems.ProblemPhyXMini0839.RadialForceDirection}
  Direction of the stirrer's net electric force along its radial axis
\end{definition}
\begin{proof}
  This is the declared model definition of Radial Force Direction.
\end{proof}
\begin{definition}[Charging Method]
  \label{def:physics:phyx-mini-0839:phyxminiproblems-problemphyxmini0839-chargingmethod}
  \lean{PhyXMiniProblems.ProblemPhyXMini0839.ChargingMethod}
  Method by which the plastic rod is charged in the prose
\end{definition}
\begin{proof}
  This is the declared model definition of Charging Method.
\end{proof}
\begin{definition}[Launch Event]
  \label{def:physics:phyx-mini-0839:phyxminiproblems-problemphyxmini0839-launchevent}
  \lean{PhyXMiniProblems.ProblemPhyXMini0839.LaunchEvent}
  Event at which the requested electric force launches the stirrer
\end{definition}
\begin{proof}
  This is the declared model definition of Launch Event.
\end{proof}
\begin{definition}[Charged Stirrer Figure]
  \label{def:physics:phyx-mini-0839:phyxminiproblems-problemphyxmini0839-chargedstirrerfigure}
  \lean{PhyXMiniProblems.ProblemPhyXMini0839.ChargedStirrerFigure}
  \uses{def:physics:phyx-mini-0839:phyxminiproblems-problemphyxmini0839-figureobject, def:physics:phyx-mini-0839:phyxminiproblems-problemphyxmini0839-figurelabel}
  Literal qualitative and scalar information read from the primary bitmap. There is deliberately no net-force value or answer-choice field here
\end{definition}
\begin{proof}
  This is the declared model definition of Charged Stirrer Figure.
\end{proof}
\begin{definition}[Charged Stirrer Setup]
  \label{def:physics:phyx-mini-0839:phyxminiproblems-problemphyxmini0839-chargedstirrersetup}
  \lean{PhyXMiniProblems.ProblemPhyXMini0839.ChargedStirrerSetup}
  \uses{def:physics:phyx-mini-0839:phyxminiproblems-problemphyxmini0839-lengthquantity, def:physics:phyx-mini-0839:phyxminiproblems-problemphyxmini0839-chargemagnitudequantity, def:physics:phyx-mini-0839:phyxminiproblems-problemphyxmini0839-linechargedensitymagnitudequantity, def:physics:phyx-mini-0839:phyxminiproblems-problemphyxmini0839-forcemagnitudequantity, def:physics:phyx-mini-0839:phyxminiproblems-problemphyxmini0839-coulombconstantquantity, def:physics:phyx-mini-0839:phyxminiproblems-problemphyxmini0839-chargesign, def:physics:phyx-mini-0839:phyxminiproblems-problemphyxmini0839-wiremodel, def:physics:phyx-mini-0839:phyxminiproblems-problemphyxmini0839-stirrerchargedistribution, def:physics:phyx-mini-0839:phyxminiproblems-problemphyxmini0839-wirestirrerorientation, def:physics:phyx-mini-0839:phyxminiproblems-problemphyxmini0839-radialforcedirection, def:physics:phyx-mini-0839:phyxminiproblems-problemphyxmini0839-chargingmethod, def:physics:phyx-mini-0839:phyxminiproblems-problemphyxmini0839-launchevent, def:physics:phyx-mini-0839:phyxminiproblems-problemphyxmini0839-chargedstirrerfigure}
  Independent physical quantities in the wire--stirrer system. In particular, the net force is not defined from the recorded answer, and the electric-field readout is not defined to be the desired final force formula
\end{definition}
\begin{proof}
  This is the declared model definition of Charged Stirrer Setup.
\end{proof}
\begin{definition}[Has Charged Stirrer Problem Data]
  \label{def:physics:phyx-mini-0839:phyxminiproblems-problemphyxmini0839-haschargedstirrerproblemdata}
  \lean{PhyXMiniProblems.ProblemPhyXMini0839.HasChargedStirrerProblemData}
  \uses{def:physics:phyx-mini-0839:phyxminiproblems-problemphyxmini0839-lengthinmeters, def:physics:phyx-mini-0839:phyxminiproblems-problemphyxmini0839-chargeincoulombs, def:physics:phyx-mini-0839:phyxminiproblems-problemphyxmini0839-linechargedensityincoulombspermeter, def:physics:phyx-mini-0839:phyxminiproblems-problemphyxmini0839-chargedstirrersetup}
  Numerical data stated in the prose, written as exact coherent-SI readouts. No force value from any answer choice occurs in this premise
\end{definition}
\begin{proof}
  This is the declared model definition of Has Charged Stirrer Problem Data.
\end{proof}
\begin{definition}[Matches Charged Stirrer Scenario]
  \label{def:physics:phyx-mini-0839:phyxminiproblems-problemphyxmini0839-matcheschargedstirrerscenario}
  \lean{PhyXMiniProblems.ProblemPhyXMini0839.MatchesChargedStirrerScenario}
  \uses{def:physics:phyx-mini-0839:phyxminiproblems-problemphyxmini0839-chargedstirrersetup}
  Qualitative modeling assumptions supplied by the prose: a long straight wire, a uniformly charged rod, like charge signs so release produces repulsion, and the radial geometry shown in the figure
\end{definition}
\begin{proof}
  This is the declared model definition of Matches Charged Stirrer Scenario.
\end{proof}
\begin{definition}[Matches Charged Stirrer Primary Figure]
  \label{def:physics:phyx-mini-0839:phyxminiproblems-problemphyxmini0839-matcheschargedstirrerprimaryfigure}
  \lean{PhyXMiniProblems.ProblemPhyXMini0839.MatchesChargedStirrerPrimaryFigure}
  \uses{def:physics:phyx-mini-0839:phyxminiproblems-problemphyxmini0839-lengthinmeters, def:physics:phyx-mini-0839:phyxminiproblems-problemphyxmini0839-linechargedensityincoulombspermeter, def:physics:phyx-mini-0839:phyxminiproblems-problemphyxmini0839-chargedstirrersetup}
  Evidence transcribed from the primary image. The three scalar labels are kept both in their printed units here and related to the setup's SI quantities
\end{definition}
\begin{proof}
  This is the declared model definition of Matches Charged Stirrer Primary Figure.
\end{proof}
\begin{definition}[Satisfies Infinite Line Charge Electrostatics]
  \label{def:physics:phyx-mini-0839:phyxminiproblems-problemphyxmini0839-satisfiesinfinitelinechargeelectrostatics}
  \lean{PhyXMiniProblems.ProblemPhyXMini0839.SatisfiesInfiniteLineChargeElectrostatics}
  \uses{def:physics:phyx-mini-0839:phyxminiproblems-problemphyxmini0839-lengthinmeters, def:physics:phyx-mini-0839:phyxminiproblems-problemphyxmini0839-chargeincoulombs, def:physics:phyx-mini-0839:phyxminiproblems-problemphyxmini0839-linechargedensityincoulombspermeter, def:physics:phyx-mini-0839:phyxminiproblems-problemphyxmini0839-forceinnewtons, def:physics:phyx-mini-0839:phyxminiproblems-problemphyxmini0839-coulombconstantinnewtonmeterssquaredpercoulombsquared, def:physics:phyx-mini-0839:phyxminiproblems-problemphyxmini0839-chargedstirrersetup}
  The electrostatic laws needed for the computation, stated at coherent-SI readout level. The first field is the standard value of Coulomb's constant. The second is the field magnitude of an infinite line charge. The third is continuous-charge superposition for a uniform radial rod: dq/dr = Q/L and dF = E dq. The last field is the general direction law for like charges. None of these fields states the requested numerical net force
\end{definition}
\begin{proof}
  This is the declared model definition of Satisfies Infinite Line Charge Electrostatics.
\end{proof}
\begin{theorem}[net Electric Force Magnitude closed Form]
  \label{thm:physics:phyx-mini-0839:phyxminiproblems-problemphyxmini0839-netelectricforcemagnitude-closedform}
  \lean{PhyXMiniProblems.ProblemPhyXMini0839.netElectricForceMagnitude_closedForm}
  \uses{def:physics:phyx-mini-0839:phyxminiproblems-problemphyxmini0839-lengthinmeters, def:physics:phyx-mini-0839:phyxminiproblems-problemphyxmini0839-chargeincoulombs, def:physics:phyx-mini-0839:phyxminiproblems-problemphyxmini0839-linechargedensityincoulombspermeter, def:physics:phyx-mini-0839:phyxminiproblems-problemphyxmini0839-forceinnewtons, def:physics:phyx-mini-0839:phyxminiproblems-problemphyxmini0839-coulombconstantinnewtonmeterssquaredpercoulombsquared, def:physics:phyx-mini-0839:phyxminiproblems-problemphyxmini0839-chargedstirrersetup, def:physics:phyx-mini-0839:phyxminiproblems-problemphyxmini0839-haschargedstirrerproblemdata, def:physics:phyx-mini-0839:phyxminiproblems-problemphyxmini0839-matcheschargedstirrerscenario, def:physics:phyx-mini-0839:phyxminiproblems-problemphyxmini0839-matcheschargedstirrerprimaryfigure, def:physics:phyx-mini-0839:phyxminiproblems-problemphyxmini0839-satisfiesinfinitelinechargeelectrostatics}
  Integration of the inverse-radius line-charge field along the stirrer gives the logarithmic net-force formula. This is a derived relation, not a field of the governing-law premise
\end{theorem}
\begin{proof}
  The conclusion follows from the governing relations and figure readouts named in the statement of net Electric Force Magnitude closed Form.
\end{proof}
\begin{definition}[Answer Choice]
  \label{def:physics:phyx-mini-0839:phyxminiproblems-problemphyxmini0839-answerchoice}
  \lean{PhyXMiniProblems.ProblemPhyXMini0839.AnswerChoice}
  The four force magnitudes printed in the multiple-choice list
\end{definition}
\begin{proof}
  This is the declared model definition of Answer Choice.
\end{proof}
\begin{definition}[answer Choice Force In Newtons]
  \label{def:physics:phyx-mini-0839:phyxminiproblems-problemphyxmini0839-answerchoiceforceinnewtons}
  \lean{PhyXMiniProblems.ProblemPhyXMini0839.answerChoiceForceInNewtons}
  \uses{def:physics:phyx-mini-0839:phyxminiproblems-problemphyxmini0839-answerchoice}
  Literal answer-choice force readout in newtons
\end{definition}
\begin{proof}
  This is the declared model definition of answer Choice Force In Newtons.
\end{proof}
\begin{definition}[choice DDisplay Tolerance In Newtons]
  \label{def:physics:phyx-mini-0839:phyxminiproblems-problemphyxmini0839-choiceddisplaytoleranceinnewtons}
  \lean{PhyXMiniProblems.ProblemPhyXMini0839.choiceDDisplayToleranceInNewtons}
  Half of the last displayed decimal place in choice D
\end{definition}
\begin{proof}
  This is the declared model definition of choice DDisplay Tolerance In Newtons.
\end{proof}
% --- Archon declaration topology end ---
% --- Archon physics formalization source end ---
